\section{La naturaleza y los límites del pre-training}

\subsection{Las ventajas únicas del pre-training}

Ilya señala las dos ventajas centrales del pre-training:
\begin{enumerate}
    \item \textbf{Un volumen de datos extraordinariamente grande}
    \item \textbf{No es necesario pensar con qué datos entrenar}: la respuesta es «everything»
\end{enumerate}

El pre-training consiste, en esencia, en un aprendizaje por compresión de «todo lo que la proyección del mundo humano deja en el texto». Pero también es muy difícil de razonar: cuando el modelo se equivoca, es difícil saber si se debe a que algún contenido está poco representado en los datos de pre-training.

\begin{warningbox}{El pre-training no tiene un análogo humano}
Ilya afirma con claridad: «I don't think there is a human analog to pre-training». Aunque algunos lo comparan con los primeros 15 años de vida de una persona o con la evolución biológica, él considera que ninguna de estas analogías es exacta. El volumen de datos del pre-training es asombroso, mientras que un humano necesita muy pocos datos para llegar a una comprensión más profunda, y sin cometer los errores elementales que sí comete el modelo.
\end{warningbox}

\subsection{El muro de datos y el siguiente paso}

Los datos de pre-training son claramente finitos. Cuando se agoten, habrá que optar por alguna forma de «souped-up pre-training» (una receta nueva, distinta de las anteriores), o pasar al RL, o explorar otras direcciones. En cualquier caso, el compute ya es muy grande.

\begin{importantbox}{De la era del scaling de vuelta a la era de la research}
\begin{itemize}
    \item 2012--2020: la era de la research: la gente probaba ideas de todo tipo para ver cuáles funcionaban
    \item 2020--2025: la era del scaling: la palabra «scaling» se llevó todo el aire de la sala, y todos hacían lo mismo
    \item 2025+: \textbf{de vuelta a la era de la research}, pero con computadoras enormes
\end{itemize}
«Is the belief really that if you just 100x the scale, everything would be transformed? I don't think that's true.»
\end{importantbox}

\subsection{Resumen de la sección}

El pre-training es una receta magnífica pero limitada. El muro de datos obliga al campo a repensar los métodos de entrenamiento. Ilya considera que estamos pasando de la era del scaling de vuelta a la era de la research: el mismo espíritu innovador que impulsa el campo, pero con recursos de cómputo órdenes de magnitud mayores.
